\documentclass[leqno]{article}
% Shared mark scheme preamble for resources.danieljsmith.org: the pack
% preamble, unchanged, then the scheme's own layer. A scheme's .tex is
% rendered from its YAML (scripts/README.md, Mark schemes), never edited.
\input{\CurrentFilePath/pack_preamble.tex}

% The page: equal side margins, so the table runs the full width of the
% header rule. The footers' right edge follows the text's.
\geometry{left=0.8in, right=0.8in}
\makeatletter
\renewcommand{\djs@footR}[1]{%
  \put(\dimexpr1in+\hoffset+\oddsidemargin+\textwidth\relax,
    -\dimexpr\paperheight-0.32in\relax){\makebox[0pt][r]{#1}}}
\makeatother

\usepackage[export]{adjustbox}
\usepackage{tabularray}
\definecolor{msGrid}{RGB}{176, 172, 164}
\definecolor{msTint}{RGB}{245, 243, 239}

% The header on every page, the front page included, and the booklet label
% under the contents-page title. \djsQuestionPageStyle stays empty.
\newcommand{\djsSchemeHeader}[1]{%
  \djsPackHeader[markscheme]{\textit{Mark scheme} - #1}%
  \renewcommand{\djsPackTitle}{#1}%
  \renewcommand{\djsBookletLabel}{Mark scheme}%
}

% The front page: the contents heading, spaced as tocloft spaces the packs'
% "Questions", then the contents table. A question's name and page number
% both link to its first page.
\newcommand{\msContentsHeading}{%
  \vskip\cftbeforetoctitleskip
  {\noindent\cfttoctitlefont\contentsname\par}%
  \vskip\cftaftertoctitleskip}
\NewTblrEnviron{mscontents}
\SetTblrInner[mscontents]{
  colspec = {Q[wd=26mm, l] X[l] Q[wd=16mm, c] Q[wd=14mm, r]},
  width = \linewidth,
  rowsep = 4pt, colsep = 6pt,
  hlines = {0.4pt, msGrid}, vlines = {0.4pt, msGrid},
  hline{2} = {2.2pt, scarlet},
  row{1} = {font=\small\bfseries},
  row{Z} = {font=\bfseries},
  column{2} = {font=\small},
}
\newcommand{\msContentsQuestion}[1]{\hyperref[q:#1]{Question #1}}
\newcommand{\msContentsPage}[1]{\hyperref[q:#1]{\pageref*{q:#1}}}
\newcommand{\msSep}{\hspace{0.9em}}

% The conventions pages' layout: headings in the documents' own bold, tables
% ruled and padded as the scheme's, roomy lists.
\newcommand{\msHead}[1]{%
  \par\vspace{18pt}%
  {\noindent\bfseries #1\par}%
  \vspace{6pt}}
\NewTblrEnviron{mstblr}
\SetTblrInner[mstblr]{rowsep = 6pt, colsep = 6pt,
  hlines = {0.4pt, msGrid}, vlines = {0.4pt, msGrid}}
\newenvironment{msList}
  {\begin{itemize}[leftmargin=1.2em, itemsep=5pt, topsep=3pt, parsep=0pt]}
  {\end{itemize}}

% A question's first page: the link target and the PDF bookmark, as
% \questionitem gives the packs.
\newcommand{\msQnum}{}
\newcommand{\msQuestion}[1]{%
  \clearpage
  \renewcommand{\msQnum}{#1}%
  \phantomsection\label{q:#1}%
  \addcontentsline{toc}{subsection}{\protect\djsTocEntry{#1}}}

% The scheme's table: Question 16mm, Scheme the rest, Marks 16mm, Notes
% 44mm; every cell ruled in msGrid; the scarlet rule under the head row. A
% long table has no caption or "continued" foot; a continuation page opens
% with a muted line above the repeated head row. The rules across a part's
% rows, the alternative's tint and the total row are the rendered table's.
\NewTblrTheme{ms}{%
  \DefTblrTemplate{firsthead}{default}{}%
  \DefTblrTemplate{middlehead,lasthead}{default}{%
    {\small\itshape\color{djsMuted}Question \msQnum\ continued\par}%
    \vspace{6pt}}%
  \DefTblrTemplate{firstfoot,middlefoot,lastfoot}{default}{}%
}
\NewTblrEnviron{msscheme}
\SetTblrOuter[msscheme]{long, theme=ms, entry=none, label=none}
\SetTblrInner[msscheme]{
  colspec = {Q[wd=16mm, l] X[l] Q[wd=16mm, c] Q[wd=44mm, l]},
  width = \linewidth,
  rowsep = 6pt, colsep = 6pt,
  rowhead = 1,
  vlines = {0.4pt, msGrid},
  hline{1} = {0.4pt, msGrid}, hline{2} = {2.2pt, scarlet},
  hline{Z} = {0.4pt, msGrid},
  column{1} = {font=\bfseries}, column{3} = {font=\bfseries},
  column{4} = {font=\small},
  row{1} = {rowsep=4pt, font=\small\bfseries},
}

% Two marks on a row: each note starts with its code.
\newcommand{\msNote}[2]{\textbf{#1}\enspace #2}
% A full-row figure: the worked solution's own picture, scaled uniformly and
% only down, to the table less the Question column. The picture is drawn
% with the line width it has in the solutions document, at its list depth
% there (the optional argument, 1 by default: a part's working), so one
% sized by \linewidth keeps its size; tabularray sets the cell's \linewidth
% to the cell's measured width, which would shrink it at every pass. The
% widths are measured once, in a solution box at the packs' text width
% (pack_preamble.tex's margins, 0.8in and 1in).
\newlength{\msSpanW}
\makeatletter
\newcommand{\ms@keepLW}[1]{%
  \expandafter\xdef\csname ms@solnLW@#1\endcsname{\the\linewidth}}
\AtBeginDocument{%
  \setlength{\msSpanW}{\dimexpr\textwidth-16mm-28pt\relax}%
  \setbox\z@\vbox{%
    \begin{minipage}{\dimexpr\paperwidth-1.8in\relax}
      \begin{djsSolution}
        \ms@keepLW{0}%
        \begin{enumerate}\item \ms@keepLW{1}%
        \begin{enumerate}\item \ms@keepLW{2}%
        \begin{enumerate}\item \ms@keepLW{3}%
        \end{enumerate}\end{enumerate}\end{enumerate}
      \end{djsSolution}
    \end{minipage}}}
\newcommand{\msFigure}[2][1]{%
  \adjustbox{max width=\msSpanW, max totalheight=0.6\textheight, valign=t}{%
    \hsize=\csname ms@solnLW@#1\endcsname
    \linewidth=\hsize
    #2}}
\makeatother
% A figure mark's feature, in the Scheme column.
\newcommand{\msFeature}[1]{\textit{#1}}
% An alternative's title row.
\newcommand{\msAlternative}[1]{\textit{Alternative:} #1}
% Drafts only: a queried mark's code, the query's why under its note, and a
% flag in its part's place.
\newcommand{\msQueried}[2]{#1~\(\to\) #2?}
\newcommand{\msQuery}[1]{{\itshape\color{djsMuted}Query: #1}}
\newcommand{\msFlag}[2]{\textbf{#1:}\enspace #2}


\djsSchemeHeader{Tangents to Polar Curves}
\djsSitePath{Core Pure}{Polar Coordinates}

\djsLastUpdated{25 October 2026}
\begin{document}
\djsSetQuestionCount{8}
\djsAddMarks{1}{6}
\djsAddMarks{2}{6}
\djsAddMarks{3}{17}
\djsAddMarks{4}{17}
\djsAddMarks{5}{10}
\djsAddMarks{6}{10}
\djsAddMarks{7}{11}
\djsAddMarks{8}{17}
\vspace*{6pt}
\djsContentsTitle
\vspace{30pt}
\msContentsHeading
\begin{mscontents}{}
Question & Marks by part & Marks & Page \\
\msContentsQuestion{1} & & 6 & \msContentsPage{1} \\
\msContentsQuestion{2} & & 6 & \msContentsPage{2} \\
\msContentsQuestion{3} & (a)\,4\msSep (b)(i)\,4\msSep (b)(ii)\,6\msSep (b)(iii)\,3 & 17 & \msContentsPage{3} \\
\msContentsQuestion{4} & (a)\,4\msSep (b)\,6\msSep (c)\,2\msSep (d)\,2\msSep (e)\,3 & 17 & \msContentsPage{4} \\
\msContentsQuestion{5} & (a)\,4\msSep (b)\,6 & 10 & \msContentsPage{5} \\
\msContentsQuestion{6} & (a)\,4\msSep (b)\,6 & 10 & \msContentsPage{6} \\
\msContentsQuestion{7} & (a)\,4\msSep (b)\,7 & 11 & \msContentsPage{7} \\
\msContentsQuestion{8} & (a)\,1\msSep (b)\,6\msSep (c)\,10 & 17 & \msContentsPage{8} \\
\SetCell[c=2]{r} Total & & 94 & \\
\end{mscontents}
\newpage
% The conventions pages, shared by every pack's mark scheme: one source of
% truth, \input after the front page, with no strand or pack logic. The keys
% on the first side, the rules on the second. The scheme layer defines
% \msHead, msList and mstblr.
\vspace*{6pt}%
{\noindent\cfttoctitlefont Reading the mark scheme\par}%
\vspace{10pt}%
These rules apply throughout, unless a note says otherwise.

\msHead{Marks}
Every mark is worth one, so \textbf{M1 A1} on a row is two marks. Only marks
coded ft follow through.\par
\vspace{10pt}%
\begin{mstblr}{
  colspec = {Q[wd=15mm, l] X[l]},
  width = \linewidth,
  hline{2} = {2.2pt, scarlet},
  column{1} = {font=\bfseries},
  row{1} = {font=\small\bfseries, rowsep=4pt},
}
Code & Meaning \\
M1 & Method: a correct method, applied to this question. A slip doesn't cost it if the method is clear. A formula must be used, not just quoted. \\
dM1 & Dependent method: as M1, but only if the method mark before it was earned, unless its note names another. \\
A1 & Accuracy: a correct answer or step. It needs the method mark it follows. \\
A1ft & As A1, following through from an earlier answer that the note names. \\
A1~AG & The last mark of a part whose answer is printed in the question: complete working, with no errors. \\
B1 & A correct value, fact or statement. It needs no method mark. \\
B1ft & As B1, following through from an earlier answer that the note names. \\
E1 & A written explanation, interpretation or conclusion. Concluding a proof by induction needs every earlier mark. \\
\end{mstblr}

\msHead{Abbreviations}
ft and AG qualify a code in the Marks column. The others appear in the
notes.\par
\vspace{10pt}%
\begin{mstblr}{
  colspec = {Q[wd=17mm, l] X[l] Q[wd=15mm, l] X[l]},
  width = \linewidth,
  column{1} = {font=\bfseries}, column{3} = {font=\bfseries},
}
ft & follow through & AG & answer given in the question \\
cao & correct answer only, with no follow-through & cso & correct solution only: no errors in the part \\
oe & or an equivalent form & awrt & any answer that rounds to this value \\
soi & seen, or implied by later working & isw & ignore working after a correct answer \\
SC & special case: credit outside the main method & their & the student's own earlier value \\
condone & a slip that doesn't cost the mark & [\;\,] & working the mark doesn't need \\
s.f. & significant figures & d.p. & decimal places \\
\end{mstblr}

\newpage
\msHead{Reading the scheme}
\begin{msList}
\item Each row is a line of working, with the marks it earns beside it.
  Where a row carries two marks, each note starts with its code.
\item \textbf{[\textit{n}]} ends each part with its total.
\item A shaded \textit{Alternative} is a second method for the same total.
  Mark an answer by the method it follows.
\item Any other valid method earns the equivalent marks, unless the question
  names a method or says \textit{hence}.
\item A sketch earns one B1 for each feature named below its figure.
\end{msList}

\msHead{Accuracy and answers}
\begin{msList}
\item Unless a question says otherwise, give answers to 3 significant
  figures, and angles in degrees to 1 decimal place. Angles in radians take
  3 significant figures.
\item A final answer with more figures is accepted if it rounds to the
  scheme's value; one with fewer is not. An answer spoiled by rounding too
  early loses its accuracy mark.
\item Where a question asks for an exact answer, a decimal loses the final
  accuracy mark. Otherwise either is accepted.
\item Wrong or missing units cost nothing unless the question asks for them.
\item A correct answer with no working earns full marks, unless the question
  needs the working: show that, prove, hence, a method it names, or an exact
  answer. A wrong answer with no working earns nothing. Calculator answers
  are treated the same way.
\item A correct answer, once reached, isn't spoiled by later working (isw).
\end{msList}

\msHead{Value of $\boldsymbol{g}$}
\begin{msList}
\item Take $g = 9.8\,\mathrm{m\,s^{-2}}$ unless a question states a value,
  which is then the only one accepted. Answers that use $g = 9.8$ may be
  given to 2 or 3 significant figures.
\item Where a question doesn't state $g$, an answer worked with $g = 10$
  throughout earns the same marks, keeping any value printed in the
  question. The worked solutions list the final answers with $g = 10$.
\end{msList}


\msQuestion{1}
\begin{msscheme}{
  hline{3-7} = {2-Z}{0.4pt, msGrid},
  hline{8} = {0.4pt, msGrid},
  row{8} = {font=\bfseries},
}
Question & Scheme & Marks & Notes \\
1 & {\(x = r\cos\theta = 3\sin 2\theta\cos\theta\)\\[4pt] \(\Rightarrow \frac{\mathrm{d}x}{\mathrm{d}\theta} = \ldots\)} & M1 & {Writes \(x\) in terms of \(\theta\) only and differentiates, by the product rule or after any rewriting. Condone sign and coefficient slips.} \\
& {\(\frac{\mathrm{d}x}{\mathrm{d}\theta} = 6\cos 2\theta\cos\theta - 3\sin 2\theta\sin\theta\)} & A1 & {Correct derivative in any form.} \\
& {\(6\cos 2\theta\cos\theta - 3\sin 2\theta\sin\theta = 0\)\\[4pt] \(\Rightarrow 6\cos\theta(3\cos^2\theta - 2) = 0\)\\[4pt] \(\Rightarrow \cos^2\theta = \ldots\)} & M1 & {Sets their \(\frac{\mathrm{d}x}{\mathrm{d}\theta} = 0\), uses identities to reach a single trigonometric ratio of \(\theta\), and solves for it. Condone slips.} \\
& {\(\cos^2\theta = \frac23\) \((\cos\theta = 0\) gives the pole\()\)} & A1 & {Or \(\sin^2\theta = \frac13\), \(\tan\theta = \frac{1}{\sqrt2}\) oe. \(\cos\theta = 0\) rejected or not pursued; A0 if it is offered as \(P\).} \\
& {\(r = 6\sin\theta\cos\theta = 6 \times \frac{1}{\sqrt3} \times \sqrt{\frac23}\)} & M1 & {Substitutes their \(\theta\) or trigonometric ratios, with \(0 < \theta < \frac{\pi}{2}\), into \(r = 3\sin 2\theta\).} \\
& {\(OP = 2\sqrt2\)} & A1 & {cso. Any exact equivalent, such as \(\sqrt8\).} \\
\nopagebreak \SetCell[c=2]{r} Total for Question 1 & & 6 & \\
\end{msscheme}

\msQuestion{2}
\begin{msscheme}{
  hline{3-6} = {2-Z}{0.4pt, msGrid},
  hline{7} = {0.4pt, msGrid},
  row{7} = {font=\bfseries},
}
Question & Scheme & Marks & Notes \\
2 & {\(y = r\sin\theta = 2\sin 3\theta\sin\theta\)\\[4pt] \(\frac{\mathrm{d}y}{\mathrm{d}\theta} = 6\cos 3\theta\sin\theta + 2\sin 3\theta\cos\theta\)} & M1 A1 & {\msNote{M1}{Differentiates \(y = 2\sin 3\theta\sin\theta\), or an equivalent such as \(\cos 2\theta - \cos 4\theta\), with respect to \(\theta\). Condone coefficient slips.}\\ \msNote{A1}{Correct \(\frac{\mathrm{d}y}{\mathrm{d}\theta}\), oe.}} \\
& {\(3\cos 3\theta\sin\theta + \sin 3\theta\cos\theta = 0\)\\[4pt] \(\Rightarrow \sin\theta\cos\theta(16\cos^2\theta - 10) = 0\)\\[4pt] \(\Rightarrow \cos^2\theta = \ldots\)} & M1 & {Sets their \(\frac{\mathrm{d}y}{\mathrm{d}\theta} = 0\) and uses multiple-angle identities to reach and solve an equation in \(\cos^2\theta\), \(\sin^2\theta\) or \(\cos 2\theta\). Condone slips.} \\
& {\(\cos^2\theta = \frac58\ \ (\sin\theta = 0\) gives the pole)} & A1 & {Or \(\sin^2\theta = \frac38\) or \(\cos 2\theta = \frac14\), with \(\theta = 0\) not taken as \(P\).} \\
& {\(\sin\theta = \frac{\sqrt6}{4}\)\\[4pt] \(\Rightarrow r = 2(3\sin\theta - 4\sin^3\theta) = \ldots\)} & M1 & {Substitutes their \(\theta\), from a valid \(\cos^2\theta\) or equivalent, into \(r = 2\sin 3\theta\), by any means.} \\
& {\(OP = \frac{3\sqrt6}{4}\)} & A1 & {Exact, oe such as \(\sqrt{\frac{27}{8}}\).} \\
\nopagebreak \SetCell[c=2]{r} Total for Question 2 & & 6 & \\
\end{msscheme}

\msQuestion{3}
\begin{msscheme}{
  hline{3-6,8-11,13-18,20-22} = {2-Z}{0.4pt, msGrid},
  hline{7,12,19,23} = {0.4pt, msGrid},
  row{23} = {font=\bfseries},
}
Question & Scheme & Marks & Notes \\
3(a) & {\(\text{Area} = \frac12\int_{-\pi}^{\pi} 4(1+\cos\theta)^2\,\mathrm{d}\theta\)} & B1 & {Correct integral with correct limits, oe such as \(2 \times \frac12\int_0^{\pi}\). Needs the \(\frac12\).} \\
& {\(\phantom{A} = 2\int_{-\pi}^{\pi}\left(\frac32 + 2\cos\theta + \frac12\cos 2\theta\right)\mathrm{d}\theta\)\\[4pt] \(\phantom{A} = 2\left[\frac32\theta + 2\sin\theta + \frac14\sin 2\theta\right]_{-\pi}^{\pi}\)} & M1 & {Expands, uses a double-angle identity for \(\cos^2\theta\) and integrates to the form \(p\theta + q\sin\theta + s\sin 2\theta\). Condone sign and coefficient slips.} \\
& {\(\phantom{A} = 2\left(\frac32\pi - \left(-\frac32\pi\right)\right)\)} & dM1 & {Substitutes their limits into their integrated expression.} \\
& {\(\phantom{A} = 6\pi\)} & A1 & {\(6\pi\), or awrt 18.8.} \\
\nopagebreak & & [4] & \\
3(b)(i) & {\(2(1+\cos\theta) = \sin\theta\)\\[4pt] \(\Rightarrow 4(1+\cos\theta)^2 = 1 - \cos^2\theta\)} & M1 & {Equates the two expressions for \(r\) and forms an equation in one trigonometric function, e.g. by squaring and using \(\sin^2\theta = 1 - \cos^2\theta\).} \\
& {\(5\cos^2\theta + 8\cos\theta + 3 = 0\)} & A1 & {oe, terms collected.} \\
& {\((5\cos\theta + 3)(\cos\theta + 1) = 0\)\\[4pt] \(\Rightarrow \cos\theta = -\frac35\) (\(\cos\theta = -1\) gives \(O\))} & dM1 & {Solves their three-term quadratic and rejects \(\cos\theta = -1\), soi by using only the other root.} \\
& {\(\sin\theta = \frac45 \Rightarrow A\left(\frac45,\ \arccos\left(-\frac35\right)\right)\)} & A1 & {\(r = \frac45\) and \(\theta = \arccos\left(-\frac35\right)\), awrt 2.21 or awrt \(126.9^\circ\).} \\
\nopagebreak & & [4] & \\
3(b)(ii) & {\(x_A = \frac45 \times \left(-\frac35\right) = -\frac{12}{25}\)} & B1 & {oe, such as \(-0.48\).} \\
& {\(2\cos\theta(1+\cos\theta) = \text{"}-\frac{12}{25}\text{"}\)} & M1 & {Sets \(r\cos\theta\) for \(C_1\) equal to their \(x\)-coordinate of \(A\).} \\
& {\(25\cos^2\theta + 25\cos\theta + 6 = 0\)} & A1 & {oe, terms collected.} \\
& {\((5\cos\theta + 3)(5\cos\theta + 2) = 0 \Rightarrow \cos\theta = \ldots\)} & dM1 & {Solves their three-term quadratic in \(\cos\theta\).} \\
& {\(\cos\theta = -\frac25\) (\(-\frac35\) is \(A\))} & A1 & {Rejection of \(-\frac35\) soi by using only \(-\frac25\).} \\
& {\(OB = 2\left(1 - \frac25\right) = \frac65\)} & A1~AG & {\(\cos\theta = -\frac25\) substituted correctly into \(r = 2(1+\cos\theta)\), reaching \(\frac65\) without error.} \\
\nopagebreak & & [6] & \\
3(b)(iii) & {\(\sin\theta_B = \frac{\sqrt{21}}{5}\)\\[4pt] \(y_A = \frac45 \times \frac45,\quad y_B = \frac65 \times \frac{\sqrt{21}}{5}\)} & M1 & {Uses \(y = r\sin\theta\) for both \(A\) and \(B\), with \(\sin\theta\) at \(B\) found from \(\cos\theta = -\frac25\). Need not be simplified.} \\
& {\(y_A = \frac{16}{25},\quad y_B = \frac{6\sqrt{21}}{25}\)} & A1 & {oe, or awrt 0.640 and 1.10.} \\
& {\(AB = \frac{6\sqrt{21} - 16}{25} = 0.460\)} & A1 & {awrt 0.460.} \\
\nopagebreak & & [3] & \\
\nopagebreak \SetCell[c=2]{r} Total for Question 3 & & 17 & \\
\end{msscheme}

\msQuestion{4}
\begin{msscheme}{
  hline{3-6,8-13,15-16,18-19,21-23} = {2-Z}{0.4pt, msGrid},
  hline{7,14,17,20,24} = {0.4pt, msGrid},
  row{24} = {font=\bfseries},
}
Question & Scheme & Marks & Notes \\
4(a) & {\(\frac12\int_0^{2\pi} a^2(1+\sin\theta)^2\,\mathrm{d}\theta\)} & B1 & {Limits spanning one full turn, e.g. \(-\frac{\pi}{2}\) to \(\frac{3\pi}{2}\), accepted. Needs the \(\frac12\).} \\
& {\(\frac{a^2}{2}\int\left(\frac32+2\sin\theta-\frac12\cos 2\theta\right)\mathrm{d}\theta\)\\[4pt] \(\phantom{A} = \frac{a^2}{2}\left[\frac32\theta-2\cos\theta-\frac14\sin 2\theta\right]\)} & M1 & {Expands, uses a double-angle identity for \(\sin^2\theta\) and integrates to the form \(p\theta+q\cos\theta+s\sin 2\theta\). Condone sign and coefficient slips.} \\
& {\(\frac{a^2}{2}\left[(3\pi-2)-(-2)\right]\)} & dM1 & {Substitutes their limits into their integrated expression and subtracts.} \\
& {\(\frac{3\pi a^2}{2}\)} & A1 & {oe.} \\
\nopagebreak & & [4] & \\
4(b) & {\(y = a(1+\sin\theta)\sin\theta\)\\[4pt] \(\frac{\mathrm{d}y}{\mathrm{d}\theta} = a\cos\theta\sin\theta\)\\[4pt] \(\phantom{\frac{\mathrm{d}y}{\mathrm{d}\theta}} + a(1+\sin\theta)\cos\theta\)} & M1 & {Differentiates \(r\sin\theta\) with their \(r\), by the product rule or after expanding. Condone sign slips.} \\
& {\(\frac{\mathrm{d}y}{\mathrm{d}\theta} = a\cos\theta(1+2\sin\theta)\)} & A1 & {oe, e.g. \(a(\cos\theta+\sin 2\theta)\).} \\
& {\(\cos\theta = 0\) or \(\sin\theta = -\frac12\)\\[4pt] \(\Rightarrow \theta = \ldots\)} & M1 & {Sets their \(\frac{\mathrm{d}y}{\mathrm{d}\theta} = 0\) and solves at least one factor for \(\theta\).} \\
& {\(\theta = \frac{\pi}{2},\ \frac{7\pi}{6},\ \frac{11\pi}{6}\)\\[4pt] (\(\frac{3\pi}{2}\) rejected: \(r = 0\))} & A1 & {These three only, or awrt 1.57, 3.67, 5.76. A0 if \(\frac{3\pi}{2}\) is included.} \\
& {\(r = a\left(1+\sin\frac{7\pi}{6}\right) = \ldots\)} & M1 & {Substitutes a value with \(\sin\theta = -\frac12\) into \(r = a(1+\sin\theta)\).} \\
& {\(A\left(\frac a2,\frac{7\pi}{6}\right)\), \(B\left(\frac a2,\frac{11\pi}{6}\right)\)} & A1 & {Both, correctly assigned to \(A\) and \(B\). Angles awrt 3.67 and 5.76 accepted.} \\
\nopagebreak & & [6] & \\
4(c) & {\(y = 2a\sin\frac{\pi}{2} = 2a\), \(y = \frac a2\sin\frac{7\pi}{6} = -\frac a4\)\\[4pt] \(WZ = 2a - \left(-\frac a4\right)\)} & M1 & {Finds \(y = r\sin\theta\) at \(\theta = \frac{\pi}{2}\) and at their \(A\) or \(B\), and finds the distance between them.} \\
& {\(WZ = \frac{9a}{4}\)} & A1~AG & {Both \(y\)-values correct, combined without error to the printed result.} \\
\nopagebreak & & [2] & \\
4(d) & {\(\frac{3\sqrt3\,a}{2}\times\frac{9a}{4}\)} & M1 & {Multiplies the given \(WX\) by \(\frac{9a}{4}\).} \\
& {\(\frac{27\sqrt3\,a^2}{8}\)} & A1 & {oe.} \\
\nopagebreak & & [2] & \\
4(e) & {\(\frac{27\sqrt3\times 8^2}{8} = 216\sqrt3\ [\approx 374.1]\)} & B1 & {\(216\sqrt3\) or awrt 374.} \\
& {\(216\sqrt3 - \frac{3\pi\times 8^2}{2} = 216\sqrt3 - 96\pi\)} & M1 & {Subtracts their area from (a), with \(a = 8\), from their rectangle area.} \\
& {\(\approx 72.5\ (\mathrm{cm}^2)\)} & A1 & {awrt 72.5, or \(216\sqrt3-96\pi\).} \\
\nopagebreak & & [3] & \\
\nopagebreak \SetCell[c=2]{r} Total for Question 4 & & 17 & \\
\end{msscheme}

\msQuestion{5}
\begin{msscheme}{
  hline{3-6,8-12} = {2-Z}{0.4pt, msGrid},
  hline{7,13} = {0.4pt, msGrid},
  row{13} = {font=\bfseries},
}
Question & Scheme & Marks & Notes \\
5(a) & {\(x = 4(1-\cos\theta)\cos\theta\)\\[4pt] \(\phantom{x} = 4(\cos\theta - \cos^2\theta) \Rightarrow \frac{\mathrm{d}x}{\mathrm{d}\theta} = \ldots\)} & M1 & {Forms \(x = r\cos\theta\) with the given \(r\) and differentiates with respect to \(\theta\), by product rule or after expanding. Condone sign slips.} \\
& {\(\frac{\mathrm{d}x}{\mathrm{d}\theta} = 4(-\sin\theta + 2\sin\theta\cos\theta)\)} & A1 & {oe; need not be factorised.} \\
& {\(4\sin\theta(2\cos\theta - 1) = 0\)\\[4pt] \(\Rightarrow \theta = 0\) (pole, rejected) or \(\cos\theta = \frac12\)} & M1 & {Sets their \(\frac{\mathrm{d}x}{\mathrm{d}\theta}\), in terms of \(\sin\theta\) and \(\cos\theta\), equal to 0, solves for \(\theta\) and rejects \(\theta = 0\) as the pole.} \\
& {\(\theta = \frac{\pi}{3},\ r = 4\left(1 - \frac12\right) = 2\)} & A1~AG & {cso, with \(r = 2\) shown by substituting \(\theta = \frac{\pi}{3}\).} \\
\nopagebreak & & [4] & \\
5(b) & {\(\text{Area} = \int_0^{\pi/3} y\,\frac{\mathrm{d}x}{\mathrm{d}\theta}\,\mathrm{d}\theta\)\\[4pt] \(\phantom{A} = 16\int_0^{\pi/3}(1-\cos\theta)\sin^2\theta\)\\[4pt] \(\phantom{A = 16\int} \times(2\cos\theta-1)\,\mathrm{d}\theta\)} & M1 A1 & {\msNote{M1}{Uses \(\int y\frac{\mathrm{d}x}{\mathrm{d}\theta}\,\mathrm{d}\theta\) with \(y = r\sin\theta\), their \(\frac{\mathrm{d}x}{\mathrm{d}\theta}\) from (a) and limits 0 and \(\frac{\pi}{3}\).}\\ \msNote{A1}{Correct integrand in \(\theta\), oe, with limits.}} \\
& {\(16\int_0^{\pi/3}(2\cos^4\theta - 3\cos^3\theta - \cos^2\theta\)\\[4pt] \(\phantom{16\int} + 3\cos\theta - 1)\,\mathrm{d}\theta = 16\left[\ldots\right]_0^{\pi/3}\)} & M1 & {Expands, uses double-angle identities for even powers and reverse chain rule or \(1-\sin^2\theta\) for odd powers, and integrates. Condone sign and coefficient slips.} \\
& {\(16\left[\sin^3\theta + \tfrac14\sin2\theta + \tfrac1{16}\sin4\theta - \tfrac34\theta\right]_0^{\pi/3}\)} & A1 & {Correct integral, oe; need not be simplified.} \\
& {\(16\left[\left(\tfrac{\sqrt3}{2}\right)^3 + \tfrac14\cdot\tfrac{\sqrt3}{2} - \tfrac1{16}\cdot\tfrac{\sqrt3}{2} - \tfrac{\pi}{4}\right] - 0\)} & dM1 & {Substitutes \(\frac{\pi}{3}\) and 0 into their integrated expression and subtracts; the zero lower value may be implied.} \\
& {\(= \tfrac{15\sqrt3}{2} - 4\pi = \tfrac12(15\sqrt3 - 8\pi)\)} & A1~AG & {cso: correct exact sine values leading without error to the printed result.} \\
\nopagebreak & & [6] & \\
\nopagebreak \SetCell[c=2]{r} Total for Question 5 & & 10 & \\
\end{msscheme}

\msQuestion{6}
\begin{msscheme}{
  hline{3-5,7-11} = {2-Z}{0.4pt, msGrid},
  hline{6,12} = {0.4pt, msGrid},
  row{12} = {font=\bfseries},
}
Question & Scheme & Marks & Notes \\
6(a) & {\(y = (1+2\cos 2\theta)\sin\theta\)\\[4pt] \(\frac{\mathrm{d}y}{\mathrm{d}\theta} = -4\sin 2\theta\sin\theta + (1+2\cos 2\theta)\cos\theta\)} & M1 A1 & {\msNote{M1}{Differentiates \(r\sin\theta\) with \(r = 1+2\cos 2\theta\), by the product rule or after expanding. Condone sign and coefficient slips.}\\ \msNote{A1}{Any correct form.}} \\
& {Sets \(\frac{\mathrm{d}y}{\mathrm{d}\theta} = 0\):\\[4pt] \(3\cos\theta(1-4\sin^2\theta) = 0 \Rightarrow \theta = \ldots\)} & M1 & {Sets their \(\frac{\mathrm{d}y}{\mathrm{d}\theta}\) to zero and solves by trigonometric identities, reaching \(\theta = \ldots\). Not for substituting \(\theta = \frac{\pi}{6}\) to verify.} \\
& {\(\sin\theta = \tfrac12 \Rightarrow \theta = \tfrac{\pi}{6}\)\\[4pt] \(r = 1+2\cos\tfrac{\pi}{3} = 2\)} & A1~AG & {Correct equation solved without error to \(\theta = \frac{\pi}{6}\) as the only solution in range, with \(r = 2\) shown by substitution.} \\
\nopagebreak & & [4] & \\
6(b) & {\(\tfrac12\int_0^{\pi/6}(1+2\cos 2\theta)^2\,\mathrm{d}\theta\)} & B1 & {May be seen within a larger expression.} \\
& {\(\tfrac12\int(3+4\cos 2\theta+2\cos 4\theta)\,\mathrm{d}\theta\)\\[4pt] \(\phantom{A} = \tfrac12\left[3\theta+2\sin 2\theta+\tfrac12\sin 4\theta\right]_0^{\pi/6}\)} & M1 & {Expands, replaces \(\cos^2 2\theta\) by \(\frac12(1\pm\cos 4\theta)\) and integrates to the form \(p\theta+q\sin 2\theta+s\sin 4\theta\). Condone sign and coefficient slips.} \\
& {\(\phantom{A} = \tfrac12\left(\tfrac{\pi}{2}+\sqrt3+\tfrac{\sqrt3}{4}\right) = \tfrac{\pi}{4}+\tfrac{5\sqrt3}{8}\)} & dM1 A1 & {\msNote{dM1}{Substitutes \(\frac{\pi}{6}\) and 0 into their integrated expression and subtracts; implied by a correct value.}} \\
& {\(P = (\sqrt3, 1)\)\\[4pt] \(R = \tfrac12(1)(3-\sqrt3)\)\\[4pt] \(\phantom{R =} - \left[\text{"}\tfrac{\pi}{4}+\tfrac{5\sqrt3}{8}\text{"} - \tfrac12(3)(2)\sin\tfrac{\pi}{6}\right]\)} & M1 & {Combines their sector area, correctly signed, with areas of triangle \(OAP\) and the triangle \(A\), \(P\), \((3,1)\), or an equivalent polygon. Triangle areas need a correct method; condone arithmetic slips.} \\
& {\(R = 3-\tfrac{\pi}{4}-\tfrac{9\sqrt3}{8}\)} & A1 & {Any exact equivalent.} \\
\nopagebreak & & [6] & \\
\nopagebreak \SetCell[c=2]{r} Total for Question 6 & & 10 & \\
\end{msscheme}

\msQuestion{7}
\begin{msscheme}{
  hline{3-6,8-14} = {2-Z}{0.4pt, msGrid},
  hline{7,15} = {0.4pt, msGrid},
  row{15} = {font=\bfseries},
}
Question & Scheme & Marks & Notes \\
7(a) & {\(x = r\cos\theta,\ y = r\sin\theta \Rightarrow \frac{\mathrm{d}x}{\mathrm{d}\theta} = \ldots,\ \frac{\mathrm{d}y}{\mathrm{d}\theta} = \ldots\)} & M1 & {Differentiates \(r\cos\theta\) or \(r\sin\theta\) by the product rule, treating \(r\) as a function of \(\theta\). Condone sign slips.} \\
& {\(\frac{\mathrm{d}x}{\mathrm{d}\theta} = \frac{\mathrm{d}r}{\mathrm{d}\theta}\cos\theta - r\sin\theta\)\\[4pt] \(\frac{\mathrm{d}y}{\mathrm{d}\theta} = \frac{\mathrm{d}r}{\mathrm{d}\theta}\sin\theta + r\cos\theta\)} & A1 & \\
& {\(\tan\psi = \frac{\mathrm{d}y}{\mathrm{d}x} = \frac{\mathrm{d}y/\mathrm{d}\theta}{\mathrm{d}x/\mathrm{d}\theta}\)} & M1 & {Uses \(\tan\psi = \frac{\mathrm{d}y}{\mathrm{d}x}\) as their \(\frac{\mathrm{d}y}{\mathrm{d}\theta}\) divided by their \(\frac{\mathrm{d}x}{\mathrm{d}\theta}\).} \\
& {\(\tan\psi = \frac{\frac{\mathrm{d}r}{\mathrm{d}\theta}\sin\theta + r\cos\theta}{\frac{\mathrm{d}r}{\mathrm{d}\theta}\cos\theta - r\sin\theta}\)} & A1~AG & {Correct derivatives substituted, leading without error to the printed result; combined steps accepted.} \\
\nopagebreak & & [4] & \\
7(b) & {\(\tan\psi = \frac{\tan\theta + \frac{r}{r'}}{1 - \tan\theta\frac{r}{r'}}\), where \(r' = \frac{\mathrm{d}r}{\mathrm{d}\theta}\)} & M1 & {Rearranges into the form \(\frac{\tan\theta + k}{1 - k\tan\theta}\), or otherwise uses the compound angle formula for tan to relate \(\psi\), \(\theta\) and \(\frac{r}{\mathrm{d}r/\mathrm{d}\theta}\). Condone sign slips.} \\
& {\(\tan\psi = \tan(\theta + \beta)\), \(\tan\beta = \frac{r}{\mathrm{d}r/\mathrm{d}\theta}\)} & A1 & {oe, such as \(\tan(\psi - \theta) = \frac{r}{\mathrm{d}r/\mathrm{d}\theta}\).} \\
& {\(\beta = \alpha \Rightarrow \frac{\mathrm{d}r}{\mathrm{d}\theta} = r\cot\alpha\)} & A1 & {Identifies the angle between \(OP\) and the tangent with \(\alpha\). Any equivalent form, such as \(\tan\alpha = \frac{r}{\mathrm{d}r/\mathrm{d}\theta}\).} \\
& {\(\int \frac{1}{r}\,\mathrm{d}r = \int \cot\alpha\,\mathrm{d}\theta\)} & M1 & {Separates their variables and integrates at least one side correctly.} \\
& {\(\ln r = (\cot\alpha)\theta + c\)} & A1 & {Must include a constant. oe, such as \(r = Ke^{(\cot\alpha)\theta}\).} \\
& {\(\theta = 0,\ r = 2 \Rightarrow c = \ldots\ [= \ln 2]\)} & M1 & {Substitutes \(\theta = 0\), \(r = 2\) into their integrated equation, which must contain a constant, to find the constant.} \\
& {\(r = 2e^{(\cot\alpha)\theta}\)} & A1~AG & {cso. Correct constant, leading without error to the printed equation.} \\
\nopagebreak & & [7] & \\
\nopagebreak \SetCell[c=2]{r} Total for Question 7 & & 11 & \\
\end{msscheme}

\msQuestion{8}
\begin{msscheme}{
  hline{3,5-8,10-17} = {2-Z}{0.4pt, msGrid},
  hline{4,9,18} = {0.4pt, msGrid},
  row{18} = {font=\bfseries},
}
Question & Scheme & Marks & Notes \\
8(a) & {\(A(6, 0)\)} & B1 & \\
\nopagebreak & & [1] & \\
8(b) & {\(y = r\sin\theta = \frac{6\sin\theta}{2-\cos\theta}\)\\[4pt] \(\frac{\mathrm{d}y}{\mathrm{d}\theta} = \frac{6\cos\theta(2-\cos\theta) - 6\sin^2\theta}{(2-\cos\theta)^2}\)} & M1 A1 & {\msNote{M1}{Differentiates \(\frac{6\sin\theta}{2-\cos\theta}\) by the quotient or product rule, or \(r\sin\theta\) by the product rule. Condone sign slips. Not for differentiating \(r\) or \(r\cos\theta\).}\\ \msNote{A1}{Any correct form; need not be simplified.}} \\
& {\(12\cos\theta - 6(\cos^2\theta + \sin^2\theta) = 0\)\\[4pt] \(\Rightarrow \cos\theta = \ldots\)} & M1 & {Sets the numerator of their quotient-rule derivative to zero, uses \(\sin^2\theta+\cos^2\theta=1\) and solves for \(\theta\).} \\
& {\(\theta = \frac{\pi}{3}\)} & A1 & {Or awrt 1.05, with no other value in \(0\leq\theta\leq\pi\).} \\
& {\(r = \frac{6}{2 - \cos\text{"}\frac{\pi}{3}\text{"}}\)\\[4pt] \(\Rightarrow B\left(4, \frac{\pi}{3}\right)\)} & M1 A1 & {\msNote{M1}{Substitutes their \(\theta\) from \(\frac{\mathrm{d}y}{\mathrm{d}\theta}=0\) into the equation of \(C\).}\\ \msNote{A1}{Or \((4, \text{awrt } 1.05)\).}} \\
\nopagebreak & & [6] & \\
8(c) & {\(2r - r\cos\theta = 6 \Rightarrow 2r = x + 6\)\\[4pt] \(\Rightarrow 4(x^2+y^2) = (x+6)^2\)} & M1 & {Uses \(x=r\cos\theta\) and \(r^2=x^2+y^2\) to reach an equation in \(x\) and \(y\) only. Condone slips in squaring.} \\
& {\(3(x-2)^2 + 4y^2 = 48\)} & A1 & {oe, such as \(3x^2+4y^2-12x-36=0\) or \(y\) as a function of \(x\).} \\
& {\(\ell_1: y = 2\sqrt3,\quad \ell_2: x = 6\)} & B1 & {Both.} \\
& {\(\text{Area} = \int_2^6 \left(2\sqrt3 - \frac{\sqrt3}{2}\sqrt{16-(x-2)^2}\right)\mathrm{d}x\)} & M1 A1 & {\msNote{M1}{Forms line minus curve from \(x=2\) to \(6\), or \(8\sqrt3\) minus \(\int_2^6 y\,\mathrm{d}x\), with \(y\) the positive root of their Cartesian equation. An equivalent integral in \(y\) accepted.}\\ \msNote{A1}{Correct integrand and limits, oe.}} \\
& {\(x - 2 = 4\sin t,\ \mathrm{d}x = 4\cos t\,\mathrm{d}t\)\\[4pt] \(\Rightarrow \int_0^{\pi/2} 16\cos^2 t\,\mathrm{d}t\)} & M1 A1 & {\msNote{M1}{Uses \(x-2=4\sin t\) or an equivalent trigonometric substitution, replacing \(\mathrm{d}x\), to reach an integral in \(t\) only. Condone coefficient slips.}\\ \msNote{A1}{Correct integrand in \(t\), any constant factor carried correctly; limits not required.}} \\
& {\(8\int_0^{\pi/2}(1+\cos 2t)\,\mathrm{d}t\)\\[4pt] \(\phantom{8} = 8\left[t + \tfrac12\sin 2t\right]_0^{\pi/2}\)} & M1 & {Uses \(\cos^2 t=\frac12(1\pm\cos 2t)\) and integrates to the form \(at+b\sin 2t\). Condone sign and coefficient slips.} \\
& {\(8\sqrt3 - \frac{\sqrt3}{2}\times 8\left(\frac{\pi}{2}\right) = \ldots\)} & dM1 & {Applies limits in \(t\) matching their substitution, or \(2\) and \(6\) after reverting to \(x\), and combines with their rectangle term.} \\
& {\(\sqrt3(8 - 2\pi)\)} & A1 & {cso. \(\sqrt3(8-2\pi)\), or \(p=8\) and \(q=-2\) stated; \(8\sqrt3-2\pi\sqrt3\) left in that form scores A0.} \\
\nopagebreak & & [10] & \\
\nopagebreak \SetCell[c=2]{r} Total for Question 8 & & 17 & \\
\end{msscheme}

\end{document}
